\section{The quotient sphere and compatible coordinates}\label{sec:setup}

We first collapse each complementary component to a point and construct
smooth compact regions exhausting the domain.  We then define
coordinates with controlled boundary oscillation and a
compatibility condition on their level sets.  The last two lemmas give
the level-set geometry and null-image criterion used to establish that
condition.

Throughout, a domain is nonempty, open, and connected.  The case
$G=\Chat$ is immediate, so we assume that the complement is nonempty.
After a M\"obius change of coordinate, we assume that
$\infty\in G$.  Fix $R>0$ such that
\[
 \{z:|z|\ge R\}\cup\{\infty\}\subset G.
\]
We use $\dd s$ and $\dd a$ for spherical length and area on the original
sphere.  We call this the \emph{physical} sphere when distinguishing it
from a quotient introduced below.  The $L^2$ norm of a one-form is its
Dirichlet energy norm; in two dimensions this norm is unchanged by a
conformal change of metric.

\subsection{Collapsing complementary components}

A \emph{continuum} is a nonempty compact connected set; it is
\emph{nondegenerate} if it has more than one point.

\begin{lemma}[The quotient sphere]\label{lem:quotient}
Collapse each component of $\Chat\setminus G$ to a point, and leave
the points of $G$ as singleton decomposition elements.  The resulting
quotient map
\[
 \pi:\Chat\longrightarrow X
\]
is a closed monotone map onto a topological sphere.  Its restriction to
$G$ is a homeomorphism onto an open subset of $X$.  The set
\[
 B=X\setminus\pi(G)
\]
is compact and totally disconnected.  Moreover, the inverse image under
$\pi$ of every continuum in $X$ is connected.
\end{lemma}

\begin{proof}
Put $F=\Chat\setminus G$.  To apply Moore's decomposition theorem,
we check that the decomposition relation is closed and that every
element is nonseparating.  Suppose first that $x_n,y_n$ belong to the
same decomposition element and converge to $x,y$.  If infinitely many
of those elements are singletons, then $x=y$.  Otherwise, after
discarding finitely many terms, let $F_n$ be the component of $F$
containing $x_n,y_n$.  The hyperspace of nonempty compact subsets of
the sphere is compact, so a subsequence of $F_n$ converges in Hausdorff
distance to a nonempty compact set $A\subset F$.  This limit is
connected and contains $x,y$.  Thus $x,y$ lie in one component of $F$,
which proves closedness of the relation.

Now let $Q$ be a component of $F$.  The connected set $G$ lies in one
component of $\Chat\setminus Q$.  Any other component $V$ would be
disjoint from $G$, and its nonempty boundary
would lie in $Q$.  The connected set $Q\cup\overline V\subset F$
would then strictly contain $Q$, a contradiction.  Hence $Q$ is
nonseparating.  Moore's theorem now identifies $X$ with a topological
sphere; see \cite{Moore1925}.  The quotient is closed and has connected
fibers.  Since $G$ is open and saturated, $\pi|_G$ is a homeomorphism
onto an open subset.

It remains to prove the inverse-image assertion and total
disconnectedness of $B$.  Let $A\subset X$ be a continuum.  If
$\pi^{-1}(A)=P\cup Q$ were a separation into disjoint nonempty compact
sets, every fiber over $A$ would lie entirely in one of $P,Q$.
Consequently $\pi(P)$ and $\pi(Q)$ would be disjoint nonempty compact
sets separating $A$, a contradiction.  A continuum contained in $B$
therefore has connected inverse image in $F$, which must lie in a
single component of $F$.
The continuum is consequently a singleton.  Finally, $B$ is the compact
image of $F$.
\end{proof}

We identify $\pi(G)$ with $G$, orient $X$ consistently with the
physical orientation on $G$, and fix a compatible metric $d_X$ on $X$.
Since a nonempty open subset of a sphere contains a nondegenerate
continuum, $B$ has empty interior.  Thus $G$ is dense in $X$.  We refer
to $B$ as the \emph{dust}.

A \emph{dust disk} is a Jordan domain $U\subset X$ whose closure
misses $\infty$ and whose boundary is a smooth physical Jordan curve
in $G$.  Its physical inverse image is exactly the bounded Jordan
domain enclosed by that same curve.  Indeed, the decomposition
elements avoid the curve and cannot cross it, so the two Jordan sides
are saturated.  We shall use the same notation for a dust disk and its
physical inverse image when no confusion can result.

The next lemma supplies small disks around dust points and smooth
contours avoiding prescribed closed sets.  These constructions will
give the compact exhaustion of $G$.

\begin{lemma}[Small disks and contours]\label{lem:dust-disks}
The following assertions hold.
\begin{enumerate}[label=\textup{(\roman*)}]
\item If $F\subset X$ is closed, every nondegenerate continuum contained
in $F\cup B$ is contained in $F$.
\item Suppose that $F\subset X$ is closed and contains $B$, and that
$\{x\}$ is a component of $F$, with $x\ne\infty$.  Every neighborhood
of $x$ contains the closure of a dust disk containing $x$ whose boundary
misses $F$.
\item Let $A$ be a compact relatively clopen subset of $B$, and let
$W\subset X\setminus\{\infty\}$ be open with $A\subset W$.
For every $\eps>0$, there are finitely many dust disks with pairwise
disjoint closures, all contained with their closures in $W$, of
$d_X$-diameter less than $\eps$, whose union contains $A$.
The disks can be required to avoid $B\setminus A$ and to meet $A$.
\item Suppose that finitely many dust disks have closures in a Jordan
domain $V\subset X\setminus\{\infty\}$, and their boundaries avoid a
closed set $F\supset B$.  Their closures can be covered by finitely many
Jordan disks with pairwise disjoint closures in $V$ and smooth
boundaries in $G\setminus F$.
\end{enumerate}
\end{lemma}

\begin{proof}
We use two facts about compact spaces.  First, a component of a compact
Hausdorff space is the intersection of its relatively clopen
neighborhoods.  Second, if $C$ is a continuum and $P$ is a proper
closed subset of $C$, every component of $P$ meets the boundary of $P$
relative to $C$.  Indeed, a component missing that compact boundary
could be separated from it by a relatively clopen subset of $P$.
That subset would be both open and closed in $C$, a contradiction.
The second fact is the boundary-bumping principle.

For \textup{(i)}, suppose a nondegenerate continuum $C\subset F\cup B$
contains $x\notin F$.  Choose $y\in C\setminus\{x\}$ and a small
closed metric neighborhood $N$ of $x$ missing both $F$ and $y$.
The component of $C\cap N$ containing $x$ meets the boundary of $N$
by boundary bumping.  It is therefore a nondegenerate subcontinuum
contained in $B$, contradicting Lemma~\ref{lem:quotient}.

For the remaining parts, we first describe a common contour
construction on the physical sphere.
If $A_0$ is a compact subset of a physical open set $O$ away from
infinity, a smooth cutoff supported in $O$ and equal to one near $A_0$
has a regular level between zero and one.  This level is a compact
smooth one-manifold, so it has finitely many components, each a smooth
Jordan curve.  The corresponding superlevel set contains $A_0$ and
has compact closure in $O$.  Taking its outer boundary curves fills
its holes.  If $O\subset\pi^{-1}(V)$ for a Jordan domain
$V\subset X\setminus\{\infty\}$, each filled outer contour still
lies in $\pi^{-1}(V)$.  Indeed,
$\pi^{-1}(X\setminus V)$ is connected by
Lemma~\ref{lem:quotient}, contains infinity, and misses the contour;
it must therefore lie on its unbounded side.

For \textup{(ii)}, take a Jordan neighborhood $V$ of $x$
whose closure is inside the prescribed neighborhood and misses
infinity.  By the first compact-space fact above, there is a
compact relatively clopen set $A_0\subset F$ with
$x\in A_0\subset V$.  Choose a physical open neighborhood of
$\pi^{-1}(A_0)$ compactly contained in
\[
 \pi^{-1}(V)\setminus\pi^{-1}(F\setminus A_0).
\]
Apply the contour construction there.  Its boundary curves miss
$\pi^{-1}(F)$.  The connected fiber $\pi^{-1}(x)$ lies
inside one of the resulting outer contours.  Its filled disk projects
to the required dust disk.  Its closure is contained in $V$ by the
connected-complement observation above.

For \textup{(iii)}, the compact metric totally disconnected space $B$
has a clopen base.  Partition $A$ into finitely many compact relatively
clopen pieces $A_1,\ldots,A_m$, each contained in a Jordan neighborhood
$V_i$ with closure in $W\setminus(B\setminus A)$ and diameter less
than $\eps$.  The compact physical preimages of these pieces are
pairwise disjoint.  They therefore have pairwise disjoint open
neighborhoods with disjoint compact closures, each lying in
$\pi^{-1}(V_i)$ and avoiding $\pi^{-1}(B\setminus A_i)$.
Make physical contours inside these neighborhoods.  The boundary
curves obtained for all pieces are mutually disjoint and miss the
whole physical complement of $G$.  Their filled disks are either
nested or have disjoint closures.  Discard the nested ones and any
disk missing $A$.  The remaining disks still cover $A$, have disjoint
closures, and each stays inside its corresponding $V_i$.

For \textup{(iv)}, let $Q$ be the union of the given closed disks.
Its boundary is contained in the union of the given boundaries and
hence misses $F$.  Use the regular-level construction to smooth a
sufficiently small physical outer neighborhood of $\pi^{-1}(Q)$.
Take the smoothing collar around the boundary of $\pi^{-1}(Q)$
disjoint from $\pi^{-1}(F)$ and keep the neighborhood's closure in
$\pi^{-1}(V)$.  There are finitely many smooth boundary
curves by the compact one-manifold argument above.  Fill the outer
contours and discard nested contours.
The connected-complement observation keeps all filled disks in $V$,
and all new boundaries remain in $G\setminus F$.
\end{proof}

We now use part \textup{(iii)} to fix the compact exhaustion.
Additional closed obstacles can be avoided by choosing $W$ disjoint
from them, and each new cover can have its closures inside the
preceding disks.  Choose cores
\[
 K_j=X\setminus\bigcup_{U\in\mathcal U_j}U,
 \qquad K_j\subset\operatorname{int}_X K_{j+1},
 \qquad \bigcup_{j\ge1}K_j=G,
\]
where each $\mathcal U_j$ is a finite family of dust disks with
pairwise disjoint closures, covers $B$, lies on the bounded side of
$|z|=R$, and has mesh tending to zero.  Here the mesh is the maximum
of the $d_X$-diameters of the disks in the cover.  Each $K_j$ is a
compact smooth bordered region in $G$ containing infinity in its
interior.  Every compact subset of $G$ lies in the interior of all
sufficiently large cores.  For $b\in B$, write
\[
 U_j(b)\in\mathcal U_j,\qquad b\in U_j(b),
 \qquad C_j(b)=\partial U_j(b).
\]
Then
\[
 \overline{U_{j+1}(b)}\subset U_j(b),
 \qquad \bigcap_j\overline{U_j(b)}=\{b\}.
\]
Smooth intermediate curves and compact physical collars can be chosen
between successive boundary curves.  We sometimes call the curves
$C_j(b)$ themselves collars; an actual annular neighborhood will always
be specified when one is needed.

\subsection{Coordinates, bands, and compatibility}

We next specify the functions to be transferred from these cores.
For each function, we record its region of definition, any additive
period, and bounds on its energy and limiting oscillation.
Compatibility will constrain how the level sets of different
coordinates can meet at the dust.

Every nonnegative sum indexed by an arbitrary set is understood as
the supremum of its finite subsums.  In particular,
\[
 \|a\|_{\ell^2(B)}^2=\sum_{s\in B}a(s)^2.
\]
A nonnegative $\ell^2(B)$ function has countable positive support.

\begin{definition}[Coordinates]\label{def:coordinates}
A coordinate $u_i$ comes with a closed set $M_i\subset X$, called its
\emph{mask}, and has one of the following forms.
\begin{enumerate}[label=\textup{(\roman*)}]
\item A bounded smooth real function on $G$, with mask $M_i=X$.
\item A coordinate with additive period one, with mask
$M_i=\overline U$ for a dust disk $U$ and singular point
$p_i\in B\cap U$.  On the cyclic universal cover of
$\overline U\setminus\{p_i\}$, it is a smooth real function over
$G\cap M_i$, smooth up to the boundary curve, satisfying
\[
 u_i(Tx)=u_i(x)+1
\]
for the positive deck transformation $T$.
\end{enumerate}
In both cases the differential is single-valued and is required to
satisfy
\[
 \|\dd u_i\|_{2,G\cap M_i}<\infty.
\]
For a nonsingular point $s\in B\cap M_i$, choose a local lift in the
periodic case and define
\[
 I_i(s)=
 \left[\liminf_{\substack{x\to s\\x\in G\cap M_i}}u_i(x),
       \limsup_{\substack{x\to s\\x\in G\cap M_i}}u_i(x)\right].
\]
We call $I_i(s)$ the \emph{band} at $s$ and require it to be a bounded
interval with length at most $a_i(s)$, where $a_i\ge0$ belongs to
$\ell^2(B)$.
For a coordinate with period, set $a_i(p_i)=1$; no real band is
required there.  Set $a_i=0$ off the mask.  A change of local sheet
translates a band by an integer.

A finite or countable well-ordered list of coordinates carries a
countable flag set $\Sigma\subset B$ containing every positive
support $\{a_i>0\}$.  Points of $B\setminus\Sigma$ are called
\emph{free}.  Every coordinate defined at a free point has a unique
trace there, interpreted modulo an integer in the periodic case.
\end{definition}

To include dust points in the level sets, write $\T=\R/\Z$ for the
circle of levels.  For $t\in\T$, let $H_i(t)\subset M_i$ consist of
the domain points with $u_i=t\pmod{1}$, the nonsingular dust points
whose bands project to sets containing $t$ in $\T$, and the singular
point, if present.  Thus the singular point belongs to every level.
We use Haar measure of total mass one on $\T$ and product Haar
measure on tuples.  Thus
\begin{equation}\label{eq:band-incidence}
 \bigl|\{t\in\T:s\in H_i(t)\}\bigr|
 \le \min\{1,a_i(s)\}\le a_i(s)
 \qquad(s\in B).
\end{equation}

Before stating compatibility, we record that each level graph
\[
 \{(s,t)\in X\times\T:s\in H_i(t)\}
\]
is compact.  At a nonsingular dust point, the lower band endpoint is
lower semicontinuous and the upper endpoint is upper semicontinuous,
as follows directly from their neighborhood-infimum and
neighborhood-supremum definitions.  The local boundedness furnished
by a finite band also bounds all nearby bands in a common lift.
These facts give closedness at nonsingular dust points.  Smoothness
gives closedness at domain points, and the definition includes every
limiting level at a singular point.  Since the mask is closed, the
graph is closed in the compact space $X\times\T$.

\begin{definition}[Compatibility]\label{def:compatibility}
A well-ordered list with its flag set satisfies \textup{(P)} if, for
each label $j$, each nonempty finite list of lower labels
$i_1,\ldots,i_d<j$ \emph{with repetitions permitted}, and almost every
$(t,t_1,\ldots,t_d)\in\T^{d+1}$, no free point of $H_j(t)$ belongs
to a nondegenerate component of
\[
 H_{i_1}(t_1)\cup\cdots\cup H_{i_d}(t_d).
\tag{P}
\]
\end{definition}

The null-set assertion uses completed product measure.  To justify
its measurability, consider the compact hyperspace of nonempty
subcontinua of $X$.  The conditions that a continuum contain a
specified point, lie in a specified finite level
union, and have diameter at least $1/n$ are closed conditions on the
continuum, point, and level tuple.  Projection over that compact
hyperspace is closed.  Taking the union over $n$ shows that belonging
to a nondegenerate component of the level union is a Borel relation
in the point and tuple.  Requiring the point to be free is Borel, and
projection over the point makes the failure relation in
Definition~\ref{def:compatibility} analytic.  It is therefore
universally measurable.  In particular, we may prove \textup{(P)} by
slicing in its level parameters; see \cite{Kechris1995}.

Increasing $\Sigma$ preserves \textup{(P)}, as does passage to a
finite ordered sublist with the same flags.  At a countable limit
stage we take the union of the lists and of their flag sets.  Every
condition in \textup{(P)} involves only finitely many coordinates,
so previously established compatibility persists at such a stage.

\subsection{The geometry of generic level unions}

Compatibility concerns nondegenerate components of finite level
unions.  We show that almost every such union has only countably many
of these components and that its subcontinua are locally connected.
The proof uses finite physical length in $G$ to control continua of
fixed positive diameter in $X$.

We write $\HH^1_{\mathrm{sph}}$ for one-dimensional Hausdorff measure
with respect to the spherical metric.

\begin{lemma}[Generic level continua]\label{lem:level-continua}
For each finite list $i_1,\ldots,i_d$, almost every tuple
$(t_1,\ldots,t_d)$ gives a compact level union
\[
 H=\bigcup_{k=1}^d H_{i_k}(t_k)
\]
with the following properties.
\begin{enumerate}[label=\textup{(\roman*)}]
\item $\HH^1_{\mathrm{sph}}(H\cap G)<\infty$.
\item For each $\eps>0$ there is a finite bound on the cardinality of
any family of pairwise disjoint subcontinua of $H$ having
$d_X$-diameter at least $\eps$.
\item $H$ has at most countably many nondegenerate components, and
every subcontinuum of $H$ is locally connected.  Any two distinct
points in a connected relatively open subset of such a subcontinuum
can be joined by a simple arc within that subset.
\end{enumerate}
In addition, if $q\ge0$ is Borel on $G$ and $m\ge0$ is in
$\ell^2(B)$, then for each label $i$,
\begin{equation}\label{eq:weighted-level-bound}
 \begin{split}
 &\int_{\T}\left(
   \int_{H_i(t)\cap G}q\,\dd\HH^1_{\mathrm{sph}}
   +\sum_{s\in H_i(t)\cap B}m(s)\right)\dd t\\
 &\hspace{15mm}\le
   \int_{G\cap M_i}q|\dd u_i|\,\dd a
   +\sum_{s\in B}m(s)a_i(s).
 \end{split}
\end{equation}
The right side is finite when $q\in L^2(G,\dd a)$.
\end{lemma}

\begin{proof}
We first prove the weighted estimate and finite length, then derive
the topological conclusions from that length bound.
Apply the coarea formula to the smooth circle-valued map
$u_i\pmod1$ on its masked domain; see \cite{Federer1969}.
Mask-boundary portions contribute zero after integration in the level
parameter: each point of that smooth boundary occurs at just one
level.  Thus the domain term in
Equation~\eqref{eq:weighted-level-bound} equals its stated area
integral.  The dust term is bounded by
Equation~\eqref{eq:band-incidence} and Tonelli's theorem; only the
countable support of $m$ is involved.  Cauchy--Schwarz bounds the
right side when $q\in L^2$.  Taking $q=1$ and $m=0$ gives finite
length for almost every level, since the spherical area is finite.
Summing over the finite list proves \textup{(i)}.

For \textup{(ii)}, only compactness of $H$ and the finite length in
\textup{(i)} are needed.  Fix $\eps>0$ and choose a core $K$ whose
complementary dust disks have
$d_X$-diameter less than $\eps$.  Every continuum $Q\subset H$ with
$\diam_X Q\ge\eps$ meets $K$: otherwise its connectedness places it
inside one of those disjoint disks.  Compactness of the physical set
$K\subset G$ and uniform continuity of $\pi$ give $r>0$ such that
every closed spherical $r$-ball centered on $K$ is contained in $G$
and has image of $d_X$-diameter less than $\eps$.  Decrease $r$ so
these balls are ordinary embedded spherical disks.

Choose $p\in Q\cap K$.  Since the closed spherical $r$-ball about
$p$ has quotient diameter less than $\eps$, $Q$ exits this ball.
By boundary bumping, the component of the intersection containing
$p$ meets the ball boundary.  Spherical distance from $p$, a
$1$-Lipschitz function, maps this subcontinuum onto $[0,r]$.  Therefore
\[
 \HH^1_{\mathrm{sph}}(Q\cap G)\ge r.
\]
For pairwise disjoint continua these measures add.  Their number is
at most $\HH^1_{\mathrm{sph}}(H\cap G)/r$, proving \textup{(ii)}.

The bound in \textup{(ii)} now gives \textup{(iii)}.
For each $n$ there are only finitely many components
of $H$ with diameter at least $1/n$, proving countability of the
nondegenerate components.  For local connectedness of every
subcontinuum, let $L\subset H$ be any subcontinuum, let
$p\in L$, and take a sufficiently small radius $t>0$.  Let $A$ be
the component through $p$ of $L\cap\overline B_X(p,t)$.  If $A$
were not a relative neighborhood of $p$ in $L$, there would be points
$p_n\to p$ in other components of this closed intersection.  After
passing to a subsequence those components are distinct: a finite
union of components not containing $p$ is compact and misses $p$.
If the closed intersection is all of $L$, it has only one component,
so this case is already excluded.  Boundary bumping now shows that
each of the distinct components meets $\{x:d_X(x,p)=t\}$.  For large
$n$ it consequently has diameter at least $t/2$, in contradiction to
\textup{(ii)}.  Hence $A$ is a connected neighborhood of $p$.
Arbitrarily small choices of $t$ prove that $L$ is locally connected.

Finally, a connected open subset $V$ of a locally connected metric
continuum is arcwise connected.  One can see path connectedness
directly by chaining connected relatively open sets of small diameter
with closures in $V$, and then recursively refining each link inside
its predecessor.  Connectedness guarantees finite chains between
specified endpoints at every refinement.  Choose the parameter
subdivisions and image diameters to tend to zero.  Completeness gives
a continuous path between the endpoints whose image remains in $V$.
The elementary loop-erasure theorem for a path in a Hausdorff space
then supplies a simple arc with the same endpoints inside its image.
These standard continuum facts are also recorded in
\cite{Whyburn1942}.
\end{proof}

We have proved that generic level continua admit simple arcs locally.
The next lemma converts bounds along such arcs into null images of
traces, which will establish compatibility.  The use of outer measure
requires no measurability of the trace function.

\begin{lemma}[A null image criterion for traces]\label{lem:trace-null}
Let $L$ be a locally connected metric continuum, let $V\subset L$
be relatively open, and let $\nu$ be a finite Borel measure on $L$.
Suppose $F\subset V$ has $\nu$-outer measure zero and a function
$v:F\to\R$ satisfies
\[
 |v(x)-v(y)|\le A\nu(\alpha)
\]
for some finite $A\ge0$, every $x,y\in F$, and every simple arc
$\alpha\subset V$ joining $x$ and $y$.  Then $v(F)$ has Lebesgue
outer measure zero.

The same conclusion holds locally when $F$ is covered by countably
many such relatively open sets and $v$ has a real branch on each.
In particular, the corresponding image in $\T$ has Haar measure
zero for a circle-valued trace with these local branches.
\end{lemma}

\begin{proof}
For $\delta>0$, regularity of the finite Borel measure on the compact
metric space $L$ gives a relatively open set $O$ with
\[
 F\subset O\subset V,\qquad \nu(O)<\delta.
\]
Local connectedness makes each component $O_k$ of $O$ relatively
open.  There are at most countably many components, because disjoint
nonempty open sets in a second-countable space are countable.
Any two points of $F\cap O_k$ can be joined by a simple arc in $O_k$.
Consequently
\[
 \diam v(F\cap O_k)\le A\nu(O_k).
\]
Cover these real images by intervals, allowing arbitrarily small
summable excess lengths.  Countable subadditivity gives
\[
 |v(F)|^*\le A\sum_k\nu(O_k)=A\nu(O)\le A\delta.
\]
Letting $\delta$ tend to zero proves the first assertion.  Apply it
separately to the countably many local branches for the second.
\end{proof}
